\newacronym{cns}{CNS}{Système Nerveux Central}
\newacronym{snc}{SNC}{Système Nerveux Central}
\newacronym{dec}{DEC}{Classification d’évaluation des drogues}
\newacronym{iacp}{IACP}{International Association of Chiefs of Polic}
\newacronym{nhtsa}{NHTSA}{Department of Transportation National Highway Traffic Safety Administration}
\newacronym{erd}{ERD}{Experts en reconnaissance de drogues}
\newacronym{inspq}{INSPQ}{Institut national de santé publique du Québec}

\newglossaryentry{cnsdepressant}
{
    name={Dépresseurs du SNC},
    description={Les dépresseurs du \acrshort{snc} « ralentissent » le cerveau et le système nerveux central d’une personne. L'alcool est le principal problème du dépresseur \acrshort{snc}.\cite{dre:hardcore}}
}
\newglossaryentry{inhalants}{
    name={Inhalants},
    description={Nommés en référence à leur principale méthode d'ingestion, les substances inhalées sont des produits chimiques respirables, notamment des solvants volatils, les gaz propulseurs ou aérosols et certains gaz anesthésiques.\cite{dre:hardcore}}
}
\newglossaryentry{dissociativeanesthetics}
{
    name={Anesthésiques dissociatifs},
    description={Médicament qui inhibe la douleur en coupant (ou en "dissociant") la perception de la douleur par le cerveau.\cite{dre:hardcore}}
}
\newglossaryentry{narcoticanalgesics}
{
    name={Analgésiques narcotiques},
    description={L'une des sept catégories de médicaments. Les Analgésiques narcotiques comprennent l'opium, les alcaloïdes naturel de l'opium (comme la morphine, codéine et thébaine ), les dérives de l'opium, les barbituriques, les stupéfiants synthétiques et de nombreux autres médicaments.\cite{dre:hardcore}}
}
\newglossaryentry{cnsstimulant}
{
    name={Stimulants du SNC},
    description={Les stimulants du \acrshort{snc} « accélèrent » l’esprit et le système nerveux central d’une personne.\cite{dre:hardcore}}
}
\newglossaryentry{hallucinogens}
{
    name={Hallucinogènes},
    description={Les hallucinogènes altèrent la capacité d’un utilisateur à percevoir la réalité en déformant les perceptions de la vue, du son, oderurs et du toucher. Elles peuvent même provoquer une « synesthésie », un phénomène par lequel une personne « mélange » ses sens.\cite{dre:hardcore}}
}
\newglossaryentry{cannabis}{
 name={Cannabis},
 description={
 Plusieurs espèces de plantes à partir desquelles sont fabriqués la marijuana et les produits connexes (exemple : cannabis Sativa et cannabis Indicia).\cite{dre:hardcore}
 }
}
\newglossaryentry{nulleffect}{
 name={effet nul},
 description={Lorsqu’une personne prend deux médicaments qui ne provoquent pas un ou plusieurs effets particuliers,
la combinaison ne provoquera pas le ou les effets.\cite{dre:hardcore}
}
}
\newglossaryentry{additiveeffect}{
 name={effet additif},
 description={Lorsqu'une personne ingère deux substances qui provoquent le ou les mêmes effets, la combinaison provoquera effet(s) amélioré(s).\cite{dre:hardcore}}
}
\newglossaryentry{antagonisticeffect}{
 name={effet antagoniste},
 description={Lorsqu’une personne ingère deux médicaments qui provoquent des effets opposés, le résultat final est imprévisible; cela dépend
de nombreux facteurs, notamment la dose, la méthode d'ingestion, la durée de l'effet et la tolérance.\cite{dre:hardcore}}
}
\newglossaryentry{dre}{
 name={DRE},
 description={est un policier formé pour reconnaître les facultés affaiblies des conducteurs sous l'influence de drogues autre que l'alcool ou en complément de celui-ci.\cite{dre:hardcore}}
}
\newglossaryentry{overlappingeffect}{
 name={effet de chevauchement},
 description={Lorsqu'un individu prend deux drogues, dont l'une provoque un effet particulier que l'autre ne provoque pas, le la combinaison provoquera l’effet.\cite{dre:hardcore}}
}
